\documentclass[12pt]{article}

\usepackage[a4paper, margin=2.5cm]{geometry}
\usepackage{amsmath}
\usepackage{amssymb}
\usepackage{amsthm}
\usepackage[colorlinks=true, linkcolor=blue, citecolor=blue, urlcolor=blue]{hyperref}
\usepackage{booktabs}
\usepackage{natbib}

\newtheorem{theorem}{Theorem}
\newtheorem{proposition}[theorem]{Proposition}
\newtheorem{definition}[theorem]{Definition}
\newtheorem{remark}[theorem]{Remark}
\newtheorem{example}[theorem]{Example}

\newcommand{\Mcl}{\mathcal{M}_{\mathrm{cl}}}
\newcommand{\Ccal}{\mathcal{C}}
\newcommand{\Fcal}{\mathcal{F}}
\newcommand{\Hcal}{\mathcal{H}}
\newcommand{\Sadm}{\mathcal{S}_{\mathrm{adm}}}
\newcommand{\Scl}{\mathcal{S}_{\mathrm{cl}}}
\newcommand{\phig}{\varphi}

\title{Phase-Coherent Closure Dynamics\\and the Emergence of Interference}
\author{%
  Lee Smart\\[4pt]
  \textit{Institute of Vibrational Field Dynamics}\\[2pt]
  \texttt{contact@vibrationalfielddynamics.org}\\[2pt]
  \texttt{@vfd\_org}%
}
\date{April 2026}

\begin{document}

\maketitle

\noindent\textbf{Status:} Preprint (not peer-reviewed)

\begin{abstract}
Paper~XIV introduced stochastic closure dynamics, producing a probabilistic layer over the closure landscape. However, that layer is dissipative and real-valued: it lacks phase, interference, and the oscillatory structure characteristic of quantum mechanics. The present paper extends the closure framework by augmenting the state space with a phase variable, introducing candidate dynamical formulations that incorporate phase evolution alongside closure descent, and demonstrating that interference arises naturally when multiple closure-compatible paths contribute to the transition amplitude between configurations.

Three candidate formulations are developed and compared: (A) a separated configuration-phase gradient flow, (B) a Hamiltonian lift (noted as a future direction), and (C) a path-integral formulation based on a complexified Onsager--Machlup action. The complexification is the paper's central ansatz: it is the minimal extension of Paper~XIV's real path functional that introduces phase using only existing closure quantities. Once the complexification ansatz and the squared-modulus probability structure are adopted, the phase-rotation frequency and amplitude are fixed by internal consistency with Paper~XIV --- no additional phase-specific parameters are introduced. A worked double-well example demonstrates that path-dependent phase accumulation produces constructive and destructive interference in the transition amplitude. Crystallisation is extended to phase-coherent selection: the realised state is the configuration that maximises both closure compatibility and phase coherence.

The paper does not claim equivalence with quantum mechanics. It establishes the first phase-bearing dynamical layer within the closure framework and the first recovery step at the level of interference structure, and identifies the conditions under which the system reduces to the dissipative gradient flow of Paper~XII in the classical limit.
\end{abstract}

%==================================================================
\section{Introduction}\label{sec:intro}
%==================================================================

Papers~XII--XIV~\citep{SmartXII, SmartXIII, SmartXIV} established a dynamical and probabilistic layer for the closure framework: deterministic gradient flow, multi-state interaction via landscape deformation, and stochastic occupation with a Boltzmann-like stationary distribution. These results provide attractor structure, transition rates, and path-integral formulation --- but they are fundamentally \textit{dissipative and real-valued}. They lack:

\begin{itemize}
    \item phase variables and oscillatory dynamics,
    \item interference between alternative paths,
    \item non-classical probability patterns arising from phase coherence.
\end{itemize}

These are the features that distinguish quantum from classical probabilistic behaviour. The present paper introduces a phase-augmented state space and explores dynamical formulations that produce interference from the closure structure.

\subsection*{Scope and epistemic status}

This paper introduces \textit{candidate formulations} for phase-coherent closure dynamics. It does not claim equivalence with quantum mechanics, does not postulate Hilbert-space axioms, and does not derive the Schr\"odinger equation. The phase structure is introduced as a minimal extension of the closure dynamics; whether it can reproduce the full structure of quantum mechanics is an open question.

The paper resolves four specific deficiencies of the Paper~XII--XIV framework:
\begin{enumerate}
    \item[(P1)] \textbf{The missing-phase problem.} Paper~XIV's stochastic dynamics is probabilistic but not quantum: it has real path weights, no phase variables, and no interference. The present paper introduces a phase-bearing amplitude consistent with the Paper~XIV stationary measure.
    \item[(P2)] \textbf{The interference problem.} Once a complexified closure path functional is adopted, interference follows structurally from the complex amplitude sum --- without importing quantum-mechanical axioms.
    \item[(P3)] \textbf{The amplitude-construction problem.} The amplitude is not arbitrary: it is fixed by consistency with the stationary distribution of Paper~XIV as $A(\psi) = \sqrt{P_{\mathrm{st}}(\psi)}$.
    \item[(P4)] \textbf{Phase-coherent crystallisation.} Crystallisation is extended from purely real stochastic settling into attractor basins to include phase stability.
\end{enumerate}
These are real advances within the VFD programme. They constitute the first genuine QM-recovery step at the level of amplitude, phase, and interference. What remains open --- unitarity, Hilbert space, and the Schr\"odinger equation --- is discussed in Section~\ref{sec:discussion}.

%==================================================================
\section{Limitations of Stochastic Closure Dynamics}\label{sec:limitations_XIV}
%==================================================================

The stochastic dynamics of Paper~XIV produces:
\begin{itemize}
    \item a stationary distribution $P_{\mathrm{st}}(\psi) \propto \exp(-2\Fcal(\psi)/\sigma^2)$,
    \item Kramers-type transition rates between attractor basins,
    \item an Onsager--Machlup path functional with real-valued weights.
\end{itemize}

However, this structure is thermodynamic rather than quantum:
\begin{itemize}
    \item the path weights $e^{-S/\sigma^2}$ are real and positive --- there is no cancellation between paths,
    \item the stationary distribution is a Boltzmann measure, not a squared amplitude,
    \item no interference arises: the probability of reaching a configuration is the sum of path probabilities, not the squared modulus of a sum of amplitudes.
\end{itemize}

To produce quantum-like behaviour, the framework must be extended to include phase.

%==================================================================
\section{Phase-Augmented State Space}\label{sec:phase}
%==================================================================

\begin{definition}[Phase-augmented state]\label{def:augmented}
The phase-augmented state space is:
\begin{equation}
    \tilde{\mathcal{S}}_{\mathrm{adm}} = \Sadm \times S^1
\end{equation}
with elements $\Psi = (\psi, \theta)$, where $\psi \in \Sadm$ is the closure configuration and $\theta \in [0, 2\pi)$ is the phase variable.
\end{definition}

Equivalently, one may write the augmented state in polar form:
\begin{equation}\label{eq:polar}
    \Psi = A(\psi)\, e^{i\theta}
\end{equation}

The amplitude $A(\psi)$ is not a free choice: once the squared-modulus probability structure $P = |\Psi|^2 = A^2$ is adopted, consistency with the stochastic dynamics of Paper~XIV fixes it. The stationary distribution derived there is $P_{\mathrm{st}}(\psi) \propto \exp(-2\Fcal(\psi)/\sigma^2)$. Consistency then requires:
\begin{equation}\label{eq:amplitude}
    A(\psi) = \sqrt{P_{\mathrm{st}}(\psi)} = \mathcal{N}\, \exp\!\left(-\frac{\Fcal(\psi)}{\sigma^2}\right)
\end{equation}
where $\mathcal{N}$ is a normalisation constant. This is the natural amplitude consistent with the squared-modulus structure and the Paper~XIV equilibrium; it is fixed once those two ingredients are accepted. On the closure manifold ($\Fcal = 0$), $A = \mathcal{N}$ is maximal; far from closure, the amplitude is exponentially suppressed.

\begin{remark}
The phase variable $\theta$ is not postulated as a quantum-mechanical phase. It is introduced as a structural degree of freedom on the closure landscape. Whether it can be identified with the quantum phase is an open question.
\end{remark}

%==================================================================
\section{Candidate Dynamical Formulations}\label{sec:formulations}
%==================================================================

We develop three candidate formulations for phase-coherent closure dynamics. The primary formalism of this paper is \textbf{Formulation~C} (path-integral phase weighting), which directly produces interference. Formulations~A and~B provide complementary differential-equation and conservative perspectives, respectively; they are presented as supporting intuition and possible future directions, not as the paper's main result.

\subsection{Formulation A: Separated configuration-phase flow (supporting)}

The simplest differential-equation extension separates the evolution of configuration and phase:
\begin{align}
    \dot{\psi} &= -\nabla\Fcal(\psi) \label{eq:formApsi} \\
    \dot{\theta} &= \Omega(\psi) \label{eq:formAtheta}
\end{align}
where $\Omega(\psi) \geq 0$ is a \textit{phase-rotation frequency} determined by the closure structure. The augmented state $\Psi(t) = A(\psi(t))\, e^{i\theta(t)}$ then evolves with dissipative amplitude dynamics and non-dissipative phase rotation.

\textbf{Determination of $\Omega$:} Within the complexified-action ansatz adopted in Formulation~C (Section~\ref{sec:formC}), the frequency is fixed: $\Omega(\psi) = \Fcal(\psi)/\sigma^2$. This is the natural choice consistent with: (i) using only the existing closure functional, (ii) vanishing phase rotation on $\Mcl$ (where $\Fcal = 0$), and (iii) the path-integral structure. An alternative, $\Omega(\psi) = \alpha\, \mathrm{tr}(H(\psi))$ (Hessian trace), governs $\Omega$ by landscape curvature rather than height and does not vanish at closure minima; it is noted as a possible extension but is not adopted here.

\textbf{Properties:}
\begin{itemize}
    \item reduces to Paper~XII gradient flow when $\Omega = 0$ or when phase is ignored,
    \item the configuration $\psi$ evolves independently of $\theta$ (one-way coupling),
    \item the phase accumulates along trajectories: $\theta(T) = \theta(0) + \int_0^T \Omega(\psi(t))\, dt$.
\end{itemize}

\textbf{Role:} Formulation~A provides the differential-equation intuition for how phase accumulates along closure trajectories. The interference structure, however, is most naturally expressed in the path-integral language of Formulation~C.

\subsection{Formulation B: Hamiltonian lift (heuristic / future direction)}

A conservative formulation would embed closure dynamics in a canonical phase space with $(\psi, \theta)$ as conjugate-like variables. This would require defining a Hamiltonian $\Hcal(\psi, \theta)$ and a proper symplectic structure on the augmented state space.

\begin{remark}
In the present paper, the canonical conjugate structure between $\psi$ and $\theta$ is not uniquely determined. A proper Hamiltonian lift would require specifying: (i) the symplectic form on $\tilde{\mathcal{S}}_{\mathrm{adm}}$, (ii) the relationship between $\theta$ and a conjugate momentum, and (iii) conditions for energy conservation. This formulation is noted as a possible conservative extension for future work but is not developed further here.
\end{remark}

\subsection{Formulation C: Path-integral phase weighting}\label{sec:formC}

The Onsager--Machlup path functional of Paper~XIV assigns a real weight $\exp(-S[\psi]/\sigma^2)$ to each path. We propose the following \textbf{minimal complexification ansatz}: extend the real action to a complex-valued action whose real part is the Onsager--Machlup action (dissipative) and whose imaginary part is the time-integrated closure functional (oscillatory). This is the simplest extension that introduces phase while remaining entirely in terms of existing closure quantities ($\Fcal$ and $\sigma$):
\begin{equation}\label{eq:complexaction}
    S_{\mathbb{C}}[\psi(\cdot)] = \underbrace{\int_0^T \|\dot{\psi} + \nabla\Fcal(\psi)\|^2\, dt}_{S[\psi]\;\text{(Onsager--Machlup)}} \;-\; i\underbrace{\int_0^T \Fcal(\psi(t))\, dt}_{\Fcal_{\mathrm{path}}[\psi]\;\text{(closure phase)}}
\end{equation}

The path amplitude is then:
\begin{equation}\label{eq:formC}
    \mathcal{A}[\psi(\cdot)] = \exp\!\left(-\frac{S_{\mathbb{C}}[\psi]}{\sigma^2}\right) = \exp\!\left(-\frac{S[\psi]}{\sigma^2}\right) \cdot \exp\!\left(\frac{i\, \Fcal_{\mathrm{path}}[\psi]}{\sigma^2}\right)
\end{equation}

Given this ansatz, the phase-rotation frequency that was a free parameter in Formulation~A is now fixed:
\begin{equation}\label{eq:omega_derived}
    \Omega(\psi) = \frac{\Fcal(\psi)}{\sigma^2}
\end{equation}

The accumulated phase along a path is:
\begin{equation}
    \Theta[\psi(\cdot)] = \int_0^T \Omega(\psi(t))\, dt = \frac{1}{\sigma^2}\int_0^T \Fcal(\psi(t))\, dt
\end{equation}

\begin{remark}[Motivation for the complexification]
The complexified action~\eqref{eq:complexaction} is chosen as the simplest extension of Paper~XIV's real path functional that: (i) produces a complex-valued amplitude, (ii) uses only the existing closure functional $\Fcal$ (no new ingredients), (iii) assigns $\Omega = 0$ on the closure manifold (where $\Fcal = 0$, so closed states carry no phase rotation), and (iv) assigns faster phase rotation to states further from closure. Other complexifications are conceivable (e.g.\ using Hessian eigenvalues or higher-order closure invariants as the imaginary part); the present choice is the most economical. The resulting structure is analogous to the relationship between Euclidean and Lorentzian path integrals: the real part damps non-classical paths, the imaginary part generates oscillation.
\end{remark}

The total amplitude for reaching configuration $\psi_f$ from $\psi_i$ is:
\begin{equation}\label{eq:totalamp}
    \mathcal{A}(\psi_i \to \psi_f) = \sum_{\text{paths}} \mathcal{A}[\psi(\cdot)]
\end{equation}

The probability is:
\begin{equation}\label{eq:probability}
    P(\psi_i \to \psi_f) = \left|\mathcal{A}(\psi_i \to \psi_f)\right|^2
\end{equation}

\textbf{Properties:}
\begin{itemize}
    \item recovers Paper~XIV when $\Theta = 0$ (no phase): $P = \sum |e^{-S/\sigma^2}|^2$ (incoherent sum),
    \item produces interference when $\Theta \neq 0$: different paths contribute with different phases,
    \item the structure is formally analogous to the Feynman path integral with $e^{iS/\hbar}$, except that the real damping factor $e^{-S/\sigma^2}$ is retained.
\end{itemize}

\subsection{Comparative evaluation}

\begin{table}[ht]
\centering
\caption{Comparison of candidate formulations.}
\label{tab:comparison}
\begin{tabular}{@{}llll@{}}
\toprule
Feature & Formulation A & Formulation B & Formulation C \\
\midrule
Phase evolution & Explicit & Canonical & Path-accumulated \\
Dissipation & Yes (modulus) & No (conservative) & Mixed \\
$\Omega$ status & From C & Undetermined & Fixed by ansatz \\
Interference & Indirect & Indirect & Direct \\
Path-integral form & Not natural & Via action & Native \\
Classical limit & $\sigma \to 0$ & Not developed & $\sigma \to 0$ (stat.\ phase) \\
Complexity & Low & Medium & Medium \\
\bottomrule
\end{tabular}
\end{table}

\begin{remark}[Assessment]
Formulation~C is the most direct route to interference: it explicitly produces path-dependent phase contributions and reduces to Paper~XIV in the phase-free limit. Formulations~A and~B provide complementary perspectives (differential-equation and Hamiltonian, respectively) and may be more suitable for specific applications. The present paper develops Formulation~C as the primary framework and uses A and B as supporting structures.
\end{remark}

%==================================================================
\section{Path Interference and Amplitude Structure}\label{sec:interference}
%==================================================================

\subsection{Interference condition}

Consider two distinct paths $\gamma_1, \gamma_2$ from $\psi_i$ to $\psi_f$, with actions $S_1, S_2$ and accumulated phases $\Theta_1, \Theta_2$. The total amplitude is:
\begin{equation}
    \mathcal{A}_{\mathrm{tot}} = e^{-S_1/\sigma^2 + i\Theta_1} + e^{-S_2/\sigma^2 + i\Theta_2}
\end{equation}

The probability is:
\begin{align}\label{eq:interference}
    P &= |\mathcal{A}_{\mathrm{tot}}|^2 \notag \\
    &= e^{-2S_1/\sigma^2} + e^{-2S_2/\sigma^2} + 2\, e^{-(S_1+S_2)/\sigma^2}\, \cos(\Theta_1 - \Theta_2)
\end{align}

The third term is the \textit{interference term}. It vanishes when $\Theta_1 = \Theta_2$ (no phase difference) and oscillates when $\Theta_1 \neq \Theta_2$.

\begin{proposition}[Interference from phase-augmented closure dynamics]\label{prop:interference}
If two paths $\gamma_1, \gamma_2$ from $\psi_i$ to $\psi_f$ have different accumulated phases $\Theta_1 \neq \Theta_2$, then:
\begin{equation}
    |\mathcal{A}_1 + \mathcal{A}_2|^2 \neq |\mathcal{A}_1|^2 + |\mathcal{A}_2|^2
\end{equation}
That is, the transition probability exhibits interference.
\end{proposition}

\begin{remark}[On the role of interference]
The interference structure identified here follows from the adoption of complex amplitudes with phase. In this sense, it is not an independent derivation of interference, but a demonstration that the closure-complexified amplitude framework reproduces the standard interference mechanism once complex structure is introduced. The non-trivial content lies in the compatibility of this structure with the closure dynamics --- in particular, the determination of $\Omega$ and $A$ from existing closure quantities --- rather than in the existence of interference itself.
\end{remark}

This is the central result of the paper: \textit{interference arises as a structural consequence of phase-augmented closure dynamics whenever distinct paths accumulate different phases}. Note that closure dynamics alone (Paper~XII) do not produce interference --- it is the phase extension that introduces the oscillatory structure.

%==================================================================
\section{Worked Example: Double-Well System}\label{sec:example}
%==================================================================

\begin{example}[Double-well interference]\label{ex:doublewell}
Consider a one-dimensional closure functional with two symmetric minima:
\begin{equation}
    \Fcal(\psi) = \lambda\, (\psi^2 - a^2)^2
\end{equation}
with minima at $\psi = \pm a$ and a barrier at $\psi = 0$ of height $\Delta\Fcal = \lambda a^4$.

In a higher-dimensional admissible state space (which is the generic case for multi-field closure configurations), two topologically distinct paths connect $\psi_i = -a$ to $\psi_f = +a$:
\begin{itemize}
    \item $\gamma_+$: passes over the barrier through $\psi = 0$ (direct crossing),
    \item $\gamma_-$: traverses a longer route through the enlarged state space, avoiding the barrier peak.
\end{itemize}
(In a strictly one-dimensional configuration space only $\gamma_+$ exists; multiple paths require at least two configuration dimensions. The example should be understood as the one-dimensional projection of a higher-dimensional system.)

With the frequency $\Omega(\psi) = \Fcal(\psi)/\sigma^2$ fixed by the complexified-action ansatz, each path accumulates a different phase:
\begin{equation}
    \Theta_\pm = \frac{1}{\sigma^2}\int_{\gamma_\pm} \Fcal(\psi(t))\, dt = \frac{\lambda}{\sigma^2}\int_{\gamma_\pm} (\psi(t)^2 - a^2)^2\, dt
\end{equation}

The direct crossing $\gamma_+$ passes through the barrier peak ($\Fcal = \lambda a^4$) and accumulates a large phase; the longer route $\gamma_-$ (if topologically available) spends more time but at lower $\Fcal$. Generically $\Theta_+ \neq \Theta_-$ whenever the time-integrated closure cost differs between paths.

The transition amplitude is:
\begin{equation}
    \mathcal{A} = e^{-S_+/\sigma^2 + i\Theta_+} + e^{-S_-/\sigma^2 + i\Theta_-}
\end{equation}

The probability exhibits interference:
\begin{equation}
    P = e^{-2S_+/\sigma^2} + e^{-2S_-/\sigma^2} + 2\, e^{-(S_+ + S_-)/\sigma^2}\, \cos(\Theta_+ - \Theta_-)
\end{equation}

When $\Theta_+ - \Theta_- = 0$: constructive interference (maximum $P$).\\
When $\Theta_+ - \Theta_- = \pi$: destructive interference (minimum $P$).\\
When $\Theta_+ - \Theta_- = \pi/2$: the interference term vanishes, recovering the incoherent sum.

No additional phase-specific parameters are introduced beyond the existing closure functional $\Fcal$ and the stochastic scale $\sigma$: $\Omega$ is fixed once the complexified-action ansatz~\eqref{eq:complexaction} is adopted, and the interference pattern is then determined entirely by the closure landscape.
\end{example}

%==================================================================
\section{Crystallisation Under Phase Coherence}\label{sec:crystal}
%==================================================================

In Paper~XIV, crystallisation was defined as the convergence of stochastic trajectories toward closure-stable attractor basins. With phase structure, crystallisation must also account for phase coherence.

\begin{definition}[Phase-coherent crystallisation]\label{def:phasecrystal}
A configuration $\Psi^* = (\psi^*, \theta^*)$ is phase-coherently crystallised if:
\begin{enumerate}
    \item $\psi^*$ is closure-stable: $\Fcal(\psi^*) = 0$,
    \item the phase $\theta^*$ is stable under the phase dynamics: $\dot{\theta} = \Omega(\psi^*)$ is constant or zero,
    \item the configuration is a local maximum of the phase-weighted probability $|\mathcal{A}|^2$ over the accessible basin.
\end{enumerate}
\end{definition}

The realised state is the configuration that maximises both closure compatibility and phase coherence. States that are closure-stable but phase-incoherent (rapidly varying $\theta$) are suppressed; states that are both closure-stable and phase-stable are selected.

%==================================================================
\section{Toward Probability Emergence}\label{sec:probability}
%==================================================================

The transition probability~\eqref{eq:probability}:
\begin{equation}
    P(\psi_i \to \psi_f) = \left|\sum_{\text{paths}} \mathcal{A}[\psi(\cdot)]\right|^2
\end{equation}
has the same squared-modulus form as the Born rule applied to a sum of quantum amplitudes, although the underlying path weights differ from the standard quantum case because a real damping factor $e^{-S/\sigma^2}$ is retained alongside the phase.

\begin{remark}[Born-like structure]
The squared-modulus form $P = |\mathcal{A}|^2$ arises as a structural consequence of the complex-valued amplitude~\eqref{eq:formC}. Within the adopted formulation, $\Omega = \Fcal/\sigma^2$ ties the phase scale directly to the closure functional, which plays an energy-like role in the closure dynamics. However, the closure path integral differs from the quantum path integral in retaining the real damping factor $e^{-S/\sigma^2}$. For the squared-modulus structure to approach the quantum Born rule, two conditions would need to hold:
\begin{itemize}
    \item the damping factor $e^{-S/\sigma^2}$ must be negligible or factorable for the relevant paths (i.e.\ the system must be near the oscillatory regime),
    \item the sum over paths must be dominated by the classical path plus Gaussian fluctuations (stationary-phase approximation).
\end{itemize}
These conditions define the regime where the closure probability structure would approach quantum Born-rule behaviour. Outside this regime, the framework predicts a modified probability structure with residual damping --- a potentially testable distinction.
\end{remark}

%==================================================================
\section{Classical Limit}\label{sec:classical}
%==================================================================

With $\Omega = \Fcal/\sigma^2$ fixed by the complexified-action ansatz, the parameter $\sigma^2$ simultaneously controls both dissipative damping and phase oscillation frequency. The classical limit is the regime where interference effects become negligible. There are two clean limits:

\begin{enumerate}
    \item \textbf{Single-path dominance ($\sigma \to 0$):} The damping factor $\exp(-S/\sigma^2)$ becomes sharply peaked around the path of minimum action $\gamma^*$. All other paths are exponentially suppressed. The sum over paths collapses:
    \begin{equation}
        \mathcal{A}(\psi_i \to \psi_f) \;\xrightarrow{\sigma \to 0}\; \exp\!\left(-\frac{S[\gamma^*]}{\sigma^2} + \frac{i\, \Fcal_{\mathrm{path}}[\gamma^*]}{\sigma^2}\right)
    \end{equation}
    Since only one path contributes, there is no multi-path interference. The probability becomes $P \propto \exp(-2S[\gamma^*]/\sigma^2)$, and the configuration dynamics is dominated by the gradient-flow solution $\dot{\psi} = -\nabla\Fcal(\psi)$ of Paper~XII.

    This is the direct analogue of the $\hbar \to 0$ stationary-phase limit in quantum mechanics, where the Feynman path integral is dominated by the classical trajectory.

    \item \textbf{Decoherence regime (environmental coupling):} If additional noise sources (beyond the $\sigma$-controlled stochastic term) introduce random phase kicks $\delta\theta$, then the ensemble average over phase fluctuations gives:
    \begin{equation}
        \langle \cos(\Theta_1 - \Theta_2 + \delta\theta) \rangle_{\delta\theta} \;\to\; 0
    \end{equation}
    The interference term averages to zero, leaving only the incoherent sum $P = \sum |\mathcal{A}_k|^2$. This recovers the real-valued path-probability structure of Paper~XIV.
\end{enumerate}

In both limits, the system reduces to previously established dynamics: gradient flow (Paper~XII) or stochastic occupation (Paper~XIV). The phase-coherent structure is an extension that operates between these limits.

%==================================================================
\section{Relation to Quantum Mechanics}\label{sec:QM}
%==================================================================

\begin{table}[ht]
\centering
\caption{Structural comparison: quantum mechanics vs phase-coherent closure dynamics.}
\label{tab:QMcomparison}
\begin{tabular}{@{}ll@{}}
\toprule
Quantum mechanics & Phase-coherent closure dynamics \\
\midrule
Hilbert state space & Augmented admissible space $\tilde{\mathcal{S}}_{\mathrm{adm}} = \Sadm \times S^1$ (not a Hilbert space) \\
Wavefunction $\psi(x)$ & Closure amplitude-phase state $\Psi = A(\psi)e^{i\theta}$ \\
Hamiltonian $\hat{H}$ & Closure functional $\Fcal$ in an energy-like role \\
Schr\"odinger equation & No established analogue; Formulations A/B are candidate directions \\
Path integral $e^{iS/\hbar}$ & $e^{(-S + i\Fcal_{\mathrm{path}})/\sigma^2}$ (Eq.~\ref{eq:formC}) \\
Interference & Demonstrated (Prop.~\ref{prop:interference}) \\
Born rule $|\psi|^2$ & $|\mathcal{A}|^2$ (same form, not derived as physical law) \\
Superposition & Multi-path amplitude sum \\
Collapse & Phase-coherent crystallisation \\
Unitarity & Not present (dissipative component retained) \\
\bottomrule
\end{tabular}
\end{table}

\textsc{Status:} The right column contains structural analogues within the phase-coherent closure framework. Interference is demonstrated; the Born-like structure is a feature of the formulation; unitarity is absent. No equivalence with quantum mechanics is claimed.

\begin{remark}[Partial formal analogy: $\sigma^2$ and $\hbar$]\label{rem:hbar}
In quantum mechanics, the Feynman path integral assigns $\exp(iS/\hbar)$ to each path, and the classical limit is $\hbar \to 0$ (stationary phase). In the closure framework, the complexified amplitude~\eqref{eq:formC} assigns $\exp((-S + i\Fcal_{\mathrm{path}})/\sigma^2)$, and the classical limit is $\sigma^2 \to 0$ (single-path dominance, Section~\ref{sec:classical}). Within the adopted formulation, $\sigma^2$ occupies part of the formal role that $\hbar$ occupies in the quantum path integral: it controls both the width of the path-probability distribution and the scale of phase oscillation. This is a structural analogy internal to the formulation, not a physical identification. In particular, the closure path integral retains a real damping factor that has no counterpart in the standard Feynman integral, so the correspondence is partial. Whether $\sigma^2$ can be physically identified with $\hbar$ --- and under what conditions the damping factor becomes negligible, completing the analogy --- is an open question.
\end{remark}

%==================================================================
\section{Discussion and Limitations}\label{sec:discussion}
%==================================================================

\subsection{What is achieved}

\textbf{Entry ansatz.} The complexified action~\eqref{eq:complexaction} is proposed as the minimal extension of Paper~XIV's real path functional that introduces phase using only existing closure quantities. This is the paper's central assumption.

\textbf{Derived within the adopted formulation:}
\begin{itemize}
    \item The amplitude $A(\psi) = \exp(-\Fcal/\sigma^2)$ is fixed by consistency with the Paper~XIV stationary distribution and the squared-modulus probability structure (Section~\ref{sec:phase}).
    \item The phase-rotation frequency $\Omega(\psi) = \Fcal(\psi)/\sigma^2$ follows immediately from the complexified action (Eq.~\ref{eq:omega_derived}).
    \item Interference is proved (Proposition~\ref{prop:interference}): whenever distinct paths accumulate different phases, the transition probability departs from the incoherent sum.
    \item The classical limit ($\sigma \to 0$) recovers single-path dominance via stationary phase; the decoherence regime recovers Paper~XIV's incoherent probabilities.
    \item Crystallisation is extended to phase-coherent selection (Definition~\ref{def:phasecrystal}).
\end{itemize}

\textbf{Structural correspondences (not physical identifications):}
\begin{itemize}
    \item The squared-modulus probability form $P = |\mathcal{A}|^2$ is structurally analogous to the Born rule.
    \item $\sigma^2$ occupies part of the formal role of $\hbar$ (Remark~\ref{rem:hbar}), but the analogy is partial: a real damping factor is retained.
\end{itemize}

\subsection{What is not achieved}

\begin{itemize}
    \item \textbf{No full QM equivalence.} The framework shares structural features with quantum mechanics but is not equivalent: unitarity is absent, the damping factor is retained, and the Born rule is a structural feature of the formulation, not derived as a physical law.
    \item \textbf{No uniqueness proof for the complexification.} The ansatz~\eqref{eq:complexaction} is the simplest complexification using existing quantities, but it is not the only possibility. Alternative imaginary terms (e.g.\ Hessian-based) are conceivable.
    \item \textbf{No Schr\"odinger equation.} No differential equation equivalent to the Schr\"odinger equation is derived.
    \item \textbf{No Hilbert-space structure.} The augmented state space is $\Sadm \times S^1$, not a Hilbert space.
    \item \textbf{No physical identification of $\sigma^2$.} The partial formal analogy with $\hbar$ does not constitute a physical identification.
    \item \textbf{The double-well example is schematic.} It demonstrates the mechanism but does not compute numerical predictions for a physical system.
\end{itemize}

\subsection{Open questions}

\begin{itemize}
    \item Can the damping factor $e^{-S/\sigma^2}$ be absorbed into a redefined path measure, leaving a purely oscillatory integral? This would be the analogue of passing from the Euclidean to the Lorentzian path integral.
    \item Is there a Wick rotation $\sigma^2 \to i\sigma^2$ that maps the damped closure dynamics to an oscillatory quantum dynamics, and if so, under what regularity conditions does it converge?
    \item Can the augmented state space $\Sadm \times S^1$ be promoted to a Hilbert space with a suitable inner product, and does the path-integral amplitude then define a unitary propagator?
    \item Are there alternative derivations of $\Omega$ (e.g.\ from Hessian eigenvalues) that remain consistent with the complexified action but give different physical predictions?
\end{itemize}

%==================================================================
\section{Conclusion}\label{sec:conclusion}
%==================================================================

This paper solves the missing-phase problem of the VFD closure programme. Papers~XII--XIV established deterministic, interacting, and stochastic dynamics on the closure landscape, but that dynamics was dissipative and real-valued: it lacked phase, interference, and the oscillatory structure characteristic of quantum mechanics.

The complexified-action ansatz~\eqref{eq:complexaction} resolves this gap. It is the simplest extension of Paper~XIV's real path functional that introduces phase using only existing closure quantities --- no additional phase-specific parameters. Once that ansatz is adopted, the main consequences follow:
\begin{itemize}
    \item the amplitude is fixed at $A(\psi) = \exp(-\Fcal/\sigma^2)$ by consistency with Paper~XIV;
    \item the phase-rotation frequency is fixed at $\Omega = \Fcal/\sigma^2$;
    \item interference arises structurally (Proposition~\ref{prop:interference});
    \item the transition probability takes the squared-modulus form $P = |\mathcal{A}|^2$;
    \item crystallisation is extended to include phase coherence.
\end{itemize}

This constitutes the first genuine QM-recovery step within the VFD framework at the level of amplitude, phase, and interference --- all from the existing closure structure once a single minimal complexification ansatz is added.

What the paper does not claim is full quantum mechanics. Unitarity is absent, the Schr\"odinger equation is not derived, Hilbert-space structure is not established, and the partial formal analogy between $\sigma^2$ and $\hbar$ is not a physical identification. These are targets for subsequent work. The present paper establishes the phase-bearing foundation on which such extensions can be built.

%==================================================================
\bibliographystyle{plainnat}
\bibliography{references}

\end{document}
